\documentclass[11pt,a4paper]{article}
\usepackage[T1]{fontenc}
\usepackage[utf8]{inputenc}
\usepackage{lmodern,microtype}
\usepackage[a4paper,margin=24mm,headheight=15pt]{geometry}
\usepackage{amsmath,amssymb,amsthm,mathtools}
\usepackage{booktabs,longtable,array,enumitem}
\usepackage{xcolor,listings,fancyhdr,xurl}
\usepackage[hidelinks]{hyperref}
\definecolor{ink}{HTML}{17364B}
\definecolor{teal}{HTML}{126D78}
\definecolor{pale}{HTML}{F1F6F7}
\definecolor{muted}{HTML}{52636D}
\hypersetup{pdftitle={Reflection and Quantitative Hankel Certificates},
 pdfauthor={Research manuscript prepared for Keio},
 pdfsubject={Proofs for exact clauses of conjectures of Sun and Sokal; priority audit}}
\pagestyle{fancy}
\fancyhf{}
\fancyhead[L]{\small\sffamily\color{muted}REFLECTION AND HANKEL CERTIFICATES}
\fancyhead[R]{\small\sffamily\color{muted}24 NOV 2026}
\fancyfoot[C]{\small\thepage}
\renewcommand{\headrulewidth}{0.3pt}
\newtheorem{theorem}{Theorem}[section]
\newtheorem{lemma}[theorem]{Lemma}
\newtheorem{corollary}[theorem]{Corollary}
\newtheorem{proposition}[theorem]{Proposition}
\theoremstyle{remark}
\newtheorem{remark}[theorem]{Remark}
\newcommand{\Fp}{\mathbb F_p}
\newcommand{\Zp}{\mathbb Z_p}
\newcommand{\LP}{LP^+}
\setlength{\parskip}{3pt}
\setlength{\emergencystretch}{2em}
\setlist{itemsep=3pt,topsep=5pt}
\allowdisplaybreaks[2]
\lstset{language=Python,basicstyle=\ttfamily\footnotesize,
 keywordstyle=\color{teal}\bfseries,commentstyle=\color{muted},
 stringstyle=\color{ink},showstringspaces=false,breaklines=true,
 columns=fullflexible,keepspaces=true,frame=single,rulecolor=\color{pale},
 backgroundcolor=\color{pale},framesep=6pt,aboveskip=9pt,belowskip=9pt}
\title{\vspace{-1cm}\color{ink}\bfseries Reflection and Quantitative\\[3pt]
 Hankel Certificates\\[8pt]\large Congruences of Sun and a finite-depth conjecture of Sokal}
\author{Research manuscript prepared for Keio}
\date{24 November 2026}
\begin{document}
\maketitle
\begin{abstract}
We prove a complementary-binomial convolution congruence modulo $p^2$
under the degree condition $ms+h<p-1$. A reflected discrete primitive
annihilates the harmonic contribution to a row derivative; this makes
the first-order $p$-adic perturbation vanish. The theorem proves all four
formulas of Z.-W.~Sun's 2019 Conjecture 36(iii), with their stated prime
ranges. It also proves three clauses of 36(ii); the remaining clause is
a direct consequence of McCarthy--Osburn's 2008 theorem and is explicitly
identified as such. Independently, we prove Sokal's 2022 Conjecture 8(c):
for $b>0$ and $k-1<\gamma<k$, the first negative $S$-fraction coefficient
of the prescribed hypergeometric logarithmic derivative is eventually
exactly the $(2k+1)$st. Finite coefficient bounds and Hankel determinant
gaps give an explicit sufficient threshold for every parameter pair and
a uniform threshold on compact boxes. The limiting fraction and Hankel
identities are classical. Detailed comparisons record substantial prior
overlap, including existing fourth-power row congruences and a stronger
known theorem about nonreal zeros. Mathematical proof and historical
priority are kept separate: this manuscript does not certify that all
six requested originality and open-status gates have been met.
\end{abstract}
\vspace{5pt}
\noindent\colorbox{pale}{\parbox{\dimexpr\linewidth-2\fboxsep\relax}{%
\small\textbf{Scope.} The proof targets are Sun 36(iii) and Sokal 8(c),
not all of either numbered conjecture. The general congruence and the
effective sign theorem are unconditional statements. This is an
internally checked research draft, not an externally refereed priority
claim or a formal proof-assistant certificate.}}
\newpage
\tableofcontents
\newpage

\section{Targets, provenance, and the research questions}
\label{sec:scope}

The two targets predate this work and are identifiable without reference
to any newly defined constant. Z.-W.~Sun's \emph{Open conjectures on
congruences} contains Conjecture 36 on printed pages 31--32
\cite{sun2019}. Sokal's paper on hypergeometric functions in the
Laguerre--P\'olya class contains Conjecture 8 in its appendix, on printed
page 8 \cite{sokal2022}. We preserve the rational parameters, signs,
prime restrictions, normalization, and quantifiers in the clauses
proved below. A sufficient threshold used to prove Sokal's eventual
sign statement is an additional conclusion, not a replacement target.

\begin{center}
\small
\begin{tabular}{@{}>{\raggedright\arraybackslash}p{0.20\linewidth}>{\raggedright\arraybackslash}p{0.33\linewidth}>{\raggedright\arraybackslash}p{0.40\linewidth}@{}}
\toprule
Original target & Exact coverage here & Deliberate boundary of the claim\\
\midrule
Sun 36(iii), 2019 & All four displayed congruences, with the source prime conditions & No assertion that these constitute a previously unsolved, first-ever resolution\\[3pt]
Sun 36(ii), 2019 & All four clauses, one using the 2008 McCarthy--Osburn theorem & The known-theorem reduction is not counted as a new result\\[3pt]
Sokal 8(c), 2022 & Every $b>0$, integer $k\ge1$, and $\gamma\in(k-1,k)$, for all sufficiently large $c$ & Parts 8(a) and 8(b) remain unproved here; known nonreal-zero results are not relabelled\\
\bottomrule
\end{tabular}
\end{center}

\subsection{How the supplied corpus was used}
The supplied repository \cite{corpus} was fixed at commit
\texttt{52c1e4ae1f82bd3fdc8b809b9a94e2e4c6ef4cb3}.
All 256 available \texttt{.tex} files were indexed and downloaded with
Git-blob checks; selected relevant arguments were then read in detail.
This was not a line-by-line validation of all 256 manuscripts. The
repository itself presents its contents as research drafts, and none
of its unverified claims is used as an established theorem here.

Entry 230, \emph{Catalan Hankel and Sun Congruences}, suggested the
complementary-row reduction of a triangular double sum and its Taylor
expansion. Entry 255, \emph{Second Minima and Supercongruences}, suggested
looking for a double vanishing mechanism. Entry 166,
\emph{Real powers of completely monotone functions have sharp derivative
threshold twelve}, supplied a thematic
connection to Stieltjes and Hankel tests. The proofs below are written
out independently and distinguish these design cues from established
external results.

\subsection{Auxiliary questions before application}
The arithmetic auxiliary question is this: when does an antisymmetric
polynomial over $\Fp$ annihilate a harmonic-weighted complete sum?
The useful answer requires double roots at consecutive reflected
endpoints. A discrete antiderivative turns the question into vanishing
of a derivative at zero. Applying it to binomial-row derivatives gives
a second power of $p$, beyond the elementary degree cancellation
modulo $p$.

The analytic auxiliary question is different: how can finitely many
coefficients certify the sign of a specified continued-fraction
coefficient on an entire parameter half-line? We bound the change in
each logarithmic-derivative coefficient, propagate the bound to Hankel
determinants, and compare it with explicit nonzero limiting products.
This closes both the existence of the extracted coefficients and their
signs. Neither an infinite continued-fraction convergence assumption
nor a numerical search for zeros is needed.

These are two concrete constructions, rather than a claim to one
historically unprecedented universal principle. In particular, both
harmonic reflection and confluent hypergeometric continued fractions
have substantial antecedents, discussed in Section~\ref{sec:priority}.

\subsection{Notation}
We use $(a)_n=a(a+1)\cdots(a+n-1)$ with $(a)_0=1$ and
$\binom zn=z(z-1)\cdots(z-n+1)/n!$. The ring of $p$-adic integers is
$\Zp$; a congruence between its elements has its usual valuation
meaning. A formal $S$-fraction means an iteratively extracted fraction
of the form $1/(1-\alpha_1t/(1-\alpha_2t/\cdots))$ at the origin.
All claims about its initial coefficients include nonvanishing of the
necessary extraction denominators.
\section{A reflection principle for complementary binomial rows}

Let $p$ be an odd prime, let $m\in\mathbb Z_{\ge2}$, and let $h\ge1$ be odd.
For $a\in\mathbb Z_p$ define
\[
 \mathcal T_{p,m,h}(a)=
 \sum_{\substack{i,j\ge0\\i+j<p}}(2i+2j+1)^h
       \frac{(a)_i^m(1-a)_j^m}{(i!)^m(j!)^m}.
\]
Every denominator in this finite sum is a $p$-adic unit.
Write $\langle x\rangle_p\in\{0,\ldots,p-1\}$ for the least
nonnegative residue of $x\in\mathbb Z_p$.

\begin{theorem}[Degree-deficit supercongruence]\label{thm:degree-deficit}
Put
\[
 s=\min\{\langle-a\rangle_p,\langle a-1\rangle_p\}.
\]
If $ms+h<p-1$, then
\[
             \mathcal T_{p,m,h}(a)\equiv0\pmod{p^2}.
\]
\end{theorem}

The extra vanishing is supplied by the harmonic cancellation in
Lemma~\ref{lem:antiderivative}. The construction proves a derivative
congruence that is needed in addition to the convolution identity.
The discrete-primitive presentation also has a direct reflection proof;
neither presentation alone certifies historical novelty.

\begin{lemma}[Reflected discrete antiderivative]\label{lem:antiderivative}
Let $P\in\mathbb F_p[t]$ have degree at most $p-2$.
Suppose $s\in\{0,\ldots,p-1\}$, $P(-s-1-t)=-P(t)$, and $P$ has a double zero
at each of $-1,\ldots,-s$ (the last condition is empty if $s=0$).
With $H_k=\sum_{j=1}^k j^{-1}\in\mathbb F_p$, one has
\[
                    \sum_{k=0}^{p-1}P(k)H_k=0.
\]
\end{lemma}

\begin{proof}
There is a unique polynomial $A\in\mathbb F_p[t]$ of degree at most
$p-1$ such that $A(0)=0$ and $A(t+1)-A(t)=P(t)$.
Indeed the finite difference takes $t^{d+1}$ to a polynomial whose
leading coefficient is $d+1$, which is nonzero for $0\le d\le p-2$.
Interchanging two finite sums gives
\[
 \sum_{k=0}^{p-1}P(k)H_k
 =\sum_{j=1}^{p-1}\frac1j\sum_{k=j}^{p-1}P(k)
 =-\sum_{j=1}^{p-1}\frac{A(j)}j=A'(0).
\]
Here $A(p)=A(0)=0$ in $\mathbb F_p$, and the last equality uses
that $A(t)/t$ has degree at most $p-2$: its sum over $\mathbb F_p$
vanishes, and its value at zero is $A'(0)$.

The zeros of $P$ give $A(-s)=A(0)=0$. The polynomial
$A(-s-t)-A(t)$ has zero finite difference and degree at most $p-1$,
so it is constant; its value at zero is zero. Thus
$A(-s-t)=A(t)$, whence $A'(-s)=-A'(0)$.
On the other hand, the double zeros give
$A'(1-j)-A'(-j)=P'(-j)=0$ for $1\le j\le s$.
Consequently $A'(-s)=A'(0)$. Since $p$ is odd, $A'(0)=0$.
\end{proof}

\begin{remark}[A second proof of the cancellation]
Let $B=p-1-s$ and $S=\sum_{k=0}^{p-1}P(k)H_k$.
The roots of $P$ restrict this sum to $0\le k\le B$.
The identities $P(B-k)=-P(k)$ and $H_{B-k}=H_{s+k}$ in $\Fp$
give
\[
 2S=-\sum_{k=0}^{B}P(k)\sum_{\ell=1}^{s}\frac1{k+\ell}.
\]
The summand extends to the polynomial
$Q(t)=P(t)\sum_{\ell=1}^s(t+\ell)^{-1}$.
For $s>0$ its degree is at most $\deg P-1$, and the double roots
ensure that it vanishes at every omitted residue; for $s=0$ it is zero.
Its complete sum over $\Fp$ is zero, so $S=0$.
This independent derivation shows explicitly the relationship with
classical harmonic reflection, including the methods in
Meng--Sun~\cite{mengsun2015}.
\end{remark}

\begin{lemma}[A complementary row and its first derivative]\label{lem:row-jet}
For $r\in\mathbb Z_{\ge0}$ let
\[
 F_r(z)=\sum_{k=0}^{p-1}(2k-z)^r(-1)^{mk}\binom zk^m
                 \in\mathbb Z_p[z].
\]
Suppose $s\in\mathbb Z_{\ge0}$ and $ms+r<p-1$, and put $B=p-1-s$.
Then $F_r(B)\equiv0\pmod p$.
If $ms+r$ is odd, then in fact
\[
                F_r(B)=0,\qquad F'_r(B)\equiv0\pmod p.
\]
\end{lemma}

\begin{proof}
Set $C_s(t)=\binom{t+s}{s}\in\mathbb F_p[t]$ and
$P_r(t)=(2t+s+1)^rC_s(t)^m$.
For $0\le k<p$, the identity
$\binom Bk\equiv(-1)^k C_s(k)\pmod p$ gives
\[
 F_r(B)\equiv\sum_{k=0}^{p-1}P_r(k)=0\pmod p,
\]
because $\deg P_r=ms+r<p-1$.
At the integer $B$, the reflection $k\mapsto B-k$ multiplies a
summand of $F_r(B)$ by $(-1)^{mB+r}=(-1)^{ms+r}$.
Thus $F_r(B)=0$ whenever $ms+r$ is odd.

For the derivative, terms with $k>B$ have zero derivative at $B$
because $m\ge2$. For $0\le k\le B$,
\[
 \left.\frac d{dz}\binom zk^m\right|_{z=B}
   =m\binom Bk^m(H_B-H_{B-k}).
\]
All harmonic numbers here have indices below $p$.
Modulo $p$, $H_B=H_s$ and $H_{B-k}=H_{s+k}$.
For $k\le B$, use
$H_{s+k}=H_k+\sum_{\ell=1}^s(k+\ell)^{-1}$.
The expression
\[
 P_r(t)\sum_{\ell=1}^s\frac1{t+\ell}
\]
is a polynomial of degree at most $ms+r-1$.
It vanishes at $B+1,\ldots,p-1$, because $m\ge2$.
Thus all these finite sums can be extended to $0\le k<p$ without
introducing undefined harmonic numbers or nonzero extra terms.
When $r=0$ the weight-derivative term is absent. When $r>0$,
that term extends to the polynomial $-r(2t+s+1)^{r-1}C_s(t)^m$.
The terms from differentiation of $(2k-z)^r$, the term $mH_sP_r$,
and the displayed polynomial all have degree below $p-1$, so their
complete sums vanish. We obtain
\[
                  F'_r(B)\equiv-m\sum_{k=0}^{p-1}P_r(k)H_k\pmod p.
\]
If $ms+r$ is odd, $P_r(-s-1-t)=-P_r(t)$, and $P_r$ has zeros of
multiplicity at least two at $-1,\ldots,-s$.
Lemma~\ref{lem:antiderivative} proves the derivative congruence.
\end{proof}

\begin{proof}[Proof of Theorem~\ref{thm:degree-deficit}]
Put $X=-a$, $Y=a-1$ and let $A=\langle X\rangle_p$,
$B=\langle Y\rangle_p$. Then $A+B=p-1$.
By symmetry we may assume $A=s\le B$.
If $i>A$ then $\binom Xi^m\in p^m\mathbb Z_p$; likewise for the
second row beyond $B$. For $i+j\ge p$, at least one of $i>A,j>B$
holds. Hence modulo $p^2$ the triangular sum can be replaced by
the full square $0\le i,j<p$.
Since $X+Y=-1$, its weight is
$(2i+2j+1)^h=((2i-X)+(2j-Y))^h$. Therefore
\[
 \mathcal T_{p,m,h}(a)
 \equiv\sum_{r=0}^h\binom hr F_r(X)F_{h-r}(Y)\pmod{p^2}.
\]
Write $X=A+p\sigma$, $Y=B+p\tau$ with $\sigma,\tau\in\mathbb Z_p$.
The values and derivatives of every $F_r$ are $p$-integral because
$k!$ is a unit for $k<p$. Taylor expansion modulo $p^2$ gives
\[
 F_r(X)F_{h-r}(Y)\equiv
 F_r(A)F_{h-r}(B)
 +p\sigma F'_r(A)F_{h-r}(B)
 +p\tau F_r(A)F'_{h-r}(B).
\]
The first product is identically zero: $A,B$ have the same parity,
and $h$ is odd, so exactly one of its factors is killed by reflection.
The coefficient of $p\sigma$ is divisible by $p$, by the first assertion of
Lemma~\ref{lem:row-jet} applied at $B$.
For the coefficient of $p\tau$, either $mA+r$ is odd and $F_r(A)=0$, or
$mA+r$ is even and $mB+h-r$ is odd. In the latter case,
Lemma~\ref{lem:row-jet} gives $F'_{h-r}(B)\equiv0\pmod p$.
Thus every summand vanishes modulo $p^2$.
\end{proof}


\begin{remark}[Two genuine limits of the general theorem]
The modulus cannot be strengthened uniformly to $p^3$.
For $a=0$, $h=1$, and every integer $m\ge2$, only the $i=0$
row survives and
\[
 \mathcal T_{p,m,1}(0)=\sum_{j=0}^{p-1}(2j+1)=p^2.
\]
The strict degree inequality also cannot simply be replaced by a weak
one. For $p=5$, $a=4$, $m=3$, and $h=1$, one has $s=1$ and
$ms+h=p-1$, but direct exact summation gives
\[
 \mathcal T_{5,3,1}(4)=-1\,440\,005\equiv20\pmod{25}.
\]
Thus the degree condition is a substantive hypothesis, not an omitted
boundary case.
\end{remark}

\section{Literal applications to Sun's Conjecture 36}
\label{sec:sun}

For comparison with Sun's original binomial notation \cite{sun2019}, set
\begin{equation}\label{eq:original-C}
 C_{p,m,h}(a)=\sum_{k=0}^{p-1}(-1)^{mk}(2k+1)^h
       \sum_{j=0}^k\binom{-a}{j}^m\binom{a-1}{k-j}^m.
\end{equation}
The identities
$\binom{-a}{j}=(-1)^j(a)_j/j!$ and
$\binom{a-1}{k-j}=(-1)^{k-j}(1-a)_{k-j}/(k-j)!$
give the exact, termwise equality
\begin{equation}\label{eq:equivalence}
                         C_{p,m,h}(a)=\mathcal T_{p,m,h}(a).
\end{equation}
There is no transformation of the target congruence: the two factors
$(-1)^{mk}$ cancel. Thus even powers in \eqref{eq:original-C} have no
alternating sign and odd powers have the source factor $(-1)^k$.

\begin{corollary}[All four assertions of Sun 36(iii)]\label{cor:sun36iii}
For primes in the stated ranges,
\begin{align*}
 C_{p,4,3}(1/5)&\equiv0\pmod{p^2}
 &&(p>11,\ p\equiv\pm1\pmod5),\\
 C_{p,4,3}(2/5)&\equiv0\pmod{p^2}
 &&(p>11,\ p\equiv\pm2\pmod5),\\
 C_{p,5,1}(1/5)&\equiv0\pmod{p^2}
 &&(p\equiv4\pmod5),\\
 C_{p,5,1}(2/5)&\equiv0\pmod{p^2}
 &&(p\equiv3\pmod5).
\end{align*}
\end{corollary}

\begin{proof}
In the first line, the smaller residue $s$ is $(p-1)/5$ when
$p\equiv1\pmod5$, and $(p-4)/5$ when $p\equiv4\pmod5$.
In the second line it is $(p-2)/5$ when $p\equiv2\pmod5$, and
$(p-3)/5$ when $p\equiv3\pmod5$.
To check every small-prime boundary, write $p=5t+r$.
The inequality $4s+3<p-1$ is equivalent to
\[
\begin{array}{c|cccc}
r&1&2&3&4\\ \hline
\text{required condition}&t>3&t>2&t>1&t>0 .
\end{array}
\]
Every prime $p>11$ in the corresponding residue class satisfies this:
the excluded small values are either composite or at most $11$.
For the third line $s=(p-4)/5$ and $5s+1=p-3$.
For the fourth line $s=(p-3)/5$ and $5s+1=p-2$.
Theorem~\ref{thm:degree-deficit} and \eqref{eq:equivalence} now apply.
The fourth line includes $p=3$, for which $s=0$; no additional
exception is required.
\end{proof}

\begin{corollary}[All four assertions of Sun 36(ii)]\label{cor:sun36ii}
For an odd prime $p$, the following statements hold:
\begin{align*}
 C_{p,4,1}(1/4)&\equiv0\pmod{p^2}&&(p\equiv3\pmod4),\\
 \sum_{k=0}^{p-1}(2k+1)\sum_{j=0}^k
    \binom{-1/2}{j}^{5}\binom{-1/2}{k-j}^{5}
       &\equiv0\pmod{p^3}&&(p\equiv3\pmod4),\\
 C_{p,6,1}(1/6)&\equiv0\pmod{p^2}&&(p\equiv5\pmod6),\\
 C_{p,6,3}(1/6)&\equiv0\pmod{p^2}&&(p\equiv5\pmod6).
\end{align*}
The second line is obtained from an already established theorem;
it is not claimed as an independent new supercongruence.
\end{corollary}

\begin{proof}
In the first line $s=(p-3)/4$, so $4s+1=p-2<p-1$.
In the last two lines $s=(p-5)/6$, so $6s+1=p-4$ and
$6s+3=p-2$. The general theorem proves these three lines.

For the second line, use McCarthy--Osburn's Theorem 1.2
\cite{mccarthyosburn2008}, which proves the relevant case of
van Hamme's (A.2) supercongruence:
\begin{equation}\label{eq:known-MO}
 \sum_{j=0}^{(p-1)/2}(4j+1)\binom{-1/2}{j}^{5}
       \equiv0\pmod{p^3}\qquad(p\equiv3\pmod4).
\end{equation}
Put $N=(p-1)/2$ and $b_j=\binom{-1/2}{j}^{5}$.
For $N<j<p$ the identity
$\binom{-1/2}{j}=(-1)^j4^{-j}\binom{2j}{j}$ implies
$b_j\in p^5\Zp$. Thus the triangle reduces modulo $p^3$ to
$0\le i,j\le N$, a square wholly inside the original triangle.
On this square,
\[
 \sum_{i,j=0}^N(2i+2j+1)b_ib_j
 =\left(\sum_{j=0}^Nb_j\right)
  \left(\sum_{j=0}^N(4j+1)b_j\right)\equiv0\pmod{p^3}.
\]
The first factor is $p$-integral and the second satisfies
\eqref{eq:known-MO}. In particular $p=3$ is included.
Notice that the second line has no $(-1)^k$ factor; inserting one
would change the source problem.
\end{proof}

\paragraph{Additional unambiguous clauses.}
The same theorem proves the first three formulas of 36(i):
\[
\begin{array}{c|c}
\text{congruence modulo }p^2&\text{prime range}\\\hline
C_{p,4,1}(1/6)\equiv0&p>3\\
C_{p,4,3}(1/6)\equiv0&p>7\\
C_{p,5,1}(1/6)\equiv0&p>7.
\end{array}
\]
Indeed $s=(p-1)/6$ or $(p-5)/6$ according as $p\equiv1$ or $5$
modulo $6$; substitution verifies the degree bound. We make no claim
to literal coverage of all of 36(i): the retrieved source rendering
of its last formula leaves a free inner index and does not display
the required inner sum. Repairing that expression would require an
editorial interpretation, which is excluded from the principal claims.
\section{An effective finite-depth form of Sokal's Conjecture 8(c)}

Let $b,c>0$, put $a=b+\gamma$, and write
\[
 f_c(z)={}_1F_2(a;b,c;z),\qquad
 \frac{bc}{a}\frac{f_c'(-t)}{f_c(-t)}
 =\cfrac{1}{1-\cfrac{\alpha_1(c)t}{1-\cfrac{\alpha_2(c)t}{1-\ddots}}}.
\]
The continued fraction is understood formally at the origin, up to any
specified depth at which its extraction denominators are nonzero.
Here $f_c'$ denotes differentiation in its argument before evaluation
at $-t$. Recall that
${}_1F_2(a;b,c;z)=\sum_{n\ge0}(a)_n z^n/((b)_n(c)_n n!)$.
This is exactly the normalization used by Sokal, \emph{When does a
hypergeometric function ${}_pF_q$ belong to the Laguerre--P\'olya class
$LP^+$?}, J. Math. Anal. Appl. 515 (2022), 126432, Appendix,
Conjecture 8(c) \cite{sokal2022}.

\begin{theorem}[Sokal 8(c), with an explicit threshold]\label{thm:sokal}
Fix an integer $k\ge1$, $b>0$, and $k-1<\gamma<k$. There is an explicitly
computable finite number $C_*(b,\gamma,k)$, defined below solely by finite
arithmetic, such that, whenever $c\ge C_*$, the first $2k+1$ coefficients
exist and satisfy
\[
 \alpha_j(c)>0\quad(1\le j\le2k),\qquad \alpha_{2k+1}(c)<0.
\]
In particular Sokal's Conjecture 8(c), with all its original quantifiers,
holds. In fact, for every fixed integer $j\ge0$,
\begin{align*}
 \lim_{c\to\infty}c\alpha_{2j+1}(c)
   &=\frac{\gamma-j}{(b+2j)(b+2j+1)},\\
 \lim_{c\to\infty}c\alpha_{2j+2}(c)
   &=\frac{b+\gamma+j+1}{(b+2j+1)(b+2j+2)}.
\end{align*}
For these limits, $\gamma$ can be any positive noninteger and each fixed
depth is considered separately. No assertion concerning the denominator
polynomials in Conjecture 8(a), or the primitive numerator polynomials in
Conjecture 8(b), is made here.
\end{theorem}

\subsection{The limiting continued fraction, derived directly}
Denote ${}_1F_1(A;B;z)$ by $M(A;B;z)$. Comparing power-series coefficients
gives
\begin{align}
 M(A;B;z)-M(A+1;B+1;z)
 &=\frac{A-B}{B(B+1)}zM(A+1;B+2;z),\label{eq:sokal-contig-one}\\
 M(A;B;z)-M(A;B+1;z)
 &=\frac{A}{B(B+1)}zM(A+1;B+2;z).\label{eq:sokal-contig-two}
\end{align}
Define
\[
 P_{2j}(t)=M(a+j;b+2j;-t),\qquad
 P_{2j+1}(t)=M(a+j+1;b+2j+1;-t),
\]
and
\[
 \beta_{2j+1}=\frac{\gamma-j}{(b+2j)(b+2j+1)},\qquad
 \beta_{2j+2}=\frac{a+j+1}{(b+2j+1)(b+2j+2)}.
\]
The two contiguous identities say precisely
\[
 P_i=P_{i+1}-\beta_{i+1}tP_{i+2}\qquad(i\ge0).
\]
Since every $P_i(0)=1$, successive division is valid in formal power
series and yields
\[
 \frac{b}{a}\frac{M'(a;b;-t)}{M(a;b;-t)}
 =\frac{P_1}{P_0}
 =\cfrac{1}{1-\cfrac{\beta_1t}{1-\cfrac{\beta_2t}{1-\ddots}}}.
\]
This is the classical confluent Gauss continued fraction; see
DLMF 13.5.1--13.5.2 \cite{dlmf}. The elementary
derivation is included to fix every sign and index.

\subsection{Finite jets and explicit perturbation bounds}
Set $N=2k+1$, and for $0\le j\le N+1$ put
\[
 U_j=\frac{(a)_j}{(b)_j j!},\qquad K_j=\frac{j(j-1)}2.
\]
The rescaled function $F_c(z)=f_c(cz)$ has coefficients
\[
 F_c(z)=\sum_{j\ge0}u_j(c)z^j,\qquad
 u_j(c)=U_jr_j(c),\qquad
 r_j(c)=\prod_{i=0}^{j-1}\left(1+\frac{i}{c}\right)^{-1}.
\]
For every $c>0$,
\begin{equation}\label{eq:sokal-jet-error}
 0\le U_j-u_j(c)\le\frac{U_jK_j}{c}.
\end{equation}
Indeed, if $v_i=i/(c+i)$, then $0\le v_i<1$ and
$1-\prod_i(1-v_i)\le\sum_i v_i\le K_j/c$.

Write $F_c'/F_c=\sum_{j\ge0}q_j(c)z^j$ and $M'/M=\sum q_j(\infty)z^j$.
Define nonnegative finite sequences recursively by
\begin{align*}
 Q_j&=(j+1)U_{j+1}+\sum_{i=1}^{j}U_iQ_{j-i},\\
 E_j&=(j+1)U_{j+1}K_{j+1}
       +\sum_{i=1}^{j}U_i\bigl(E_{j-i}+K_iQ_{j-i}\bigr)
       \qquad(0\le j\le N).
\end{align*}
Empty sums are zero, so $Q_0=U_1$ and $E_0=0$.
From
\[
 q_j(c)=(j+1)u_{j+1}(c)-\sum_{i=1}^{j}u_i(c)q_{j-i}(c)
\]
and \eqref{eq:sokal-jet-error}, induction gives
\begin{equation}\label{eq:sokal-log-jet-error}
 |q_j(c)|,\ |q_j(\infty)|\le Q_j,\qquad
 |q_j(c)-q_j(\infty)|\le E_j/c.
\end{equation}
This estimate uses only finite jets; no uniform convergence of an infinite
continued fraction is being assumed.

Let
\[
 G_c(t)=\frac ba\frac{F_c'(-t)}{F_c(-t)}=\sum_{j\ge0}m_j(c)t^j,
 \qquad m_j(c)=\frac ba(-1)^jq_j(c),
\]
and use $G_\infty$ for the analogous Kummer quotient. In particular
$m_0(c)=m_0(\infty)=1$. Put
\[
 B=\max\left(1,\max_{0\le j\le N}\frac baQ_j\right),\qquad
 E=\max_{0\le j\le N}\frac baE_j.
\]
Then all relevant moments have absolute value at most $B$, and each
moment changes by at most $E/c$.

\subsection{The determinant sign certificate}
For $s\ge1$ and $\epsilon\in\{0,1\}$ define
\[
 H_s^{(\epsilon)}(c)
   =\det\bigl(m_{i+j+\epsilon}(c)\bigr)_{i,j=0}^{s-1},\qquad
 H_0^{(\epsilon)}=1.
\]
All matrices below use moments of index at most $N$. The standard formal
Hankel identities for an $S$-fraction give
\begin{align*}
 h_s^{(0)}:=H_s^{(0)}(\infty)
 &=\prod_{j=1}^{s-1}(\beta_{2j-1}\beta_{2j})^{s-j},\\
 h_s^{(1)}:=H_s^{(1)}(\infty)
 &=\left(\prod_{j=1}^{s}\beta_{2j-1}\right)h_s^{(0)}.
\end{align*}
These are classical algebraic identities and require no positivity
assumption; see, for example, Sokal's Theorem 9.5(b,c)
\cite{sokal2023}.
For example, contraction to a $J$-fraction gives diagonal coefficients
$\beta_1,\beta_2+\beta_3,\beta_4+\beta_5,\ldots$ and lower coefficients
$\beta_1\beta_2,\beta_3\beta_4,\ldots$; Gaussian elimination gives the first
formula. The monic polynomial recurrence at zero gives
$\pi_s(0)=(-1)^s\beta_1\beta_3\cdots\beta_{2s-1}$ for the monic
orthogonal polynomials $\pi_s$, and the usual bordered
Hankel formula gives the second.

Here $\beta_1,\ldots,\beta_{2k}>0$ and $\beta_{2k+1}<0$. Consequently
\[
 h_s^{(0)}>0\ (1\le s\le k+1),\qquad
 h_s^{(1)}>0\ (1\le s\le k),\qquad h_{k+1}^{(1)}<0.
\]
The determinant expansion and telescoping each product imply
\begin{equation}\label{eq:sokal-det-error}
 \left|H_s^{(\epsilon)}(c)-h_s^{(\epsilon)}\right|
 \le \frac{s!\,s\,E B^{s-1}}{c}.
\end{equation}
Indeed there are $s!$ permutation products, each has $s$ factors bounded
by $B$ at both endpoints, and changing one factor costs at most $E/c$.
Thus the following entirely explicit number suffices:
\begin{equation}\label{eq:sokal-effective-threshold}
 C_*(b,\gamma,k)=\max\left\{1,
 \max_{\substack{1\le s\le k+1\\\epsilon\in\{0,1\}}}
 \frac{2s!\,s\,E B^{s-1}}{|h_s^{(\epsilon)}|}\right\}.
\end{equation}
Every denominator in this expression is nonzero. For $c\ge C_*$, every
$H_s^{(\epsilon)}(c)$ has the sign of $h_s^{(\epsilon)}$ and absolute value
at least $|h_s^{(\epsilon)}|/2$.

The extraction formulas are
\begin{align*}
 \widetilde\alpha_{2s-1}(c)
 &=\frac{H_s^{(1)}(c)H_{s-1}^{(0)}(c)}
        {H_s^{(0)}(c)H_{s-1}^{(1)}(c)},\\
 \widetilde\alpha_{2s}(c)
 &=\frac{H_{s+1}^{(0)}(c)H_{s-1}^{(1)}(c)}
        {H_s^{(1)}(c)H_s^{(0)}(c)}.
\end{align*}
All needed denominators are nonzero; these formulas prove positivity up
to index $2k$ and negativity at index $2k+1$. Finally
$G_c(t)=(bc/a)f_c'(-ct)/f_c(-ct)$, so
$\widetilde\alpha_j(c)=c\alpha_j(c)$. As $c>0$, the desired signs follow.
The same formulas, together with the coefficient limits and nonzero
limiting determinants, prove all the asserted fixed-depth limits.

\subsection{A uniform explicit threshold on parameter boxes}
Fix $0<b_-\le b_+$ and $0<\eta\le1/2$, and allow
\[
 b\in[b_-,b_+],\qquad \gamma\in[k-1+\eta,k-\eta].
\]
An explicit uniform version follows without taking an unknown supremum.
Replace $U_j$ above by
\[
 U_j^+=\frac{(b_++k)_j}{(b_-)_j j!}
\]
and form $Q_j^+,E_j^+$ by the same nonnegative recursions. Put
\[
 B^+=\max(1,Q_0^+,\ldots,Q_N^+),\qquad
 E^+=\max(E_0^+,\ldots,E_N^+),\qquad
 \delta=\frac{\eta}{(b_++2k)(b_++2k+1)}.
\]
Then $0<\delta<1$ and $|\beta_j|\ge\delta$ for $1\le j\le2k+1$.
(The odd numerators have magnitude at least $\eta$; the even numerators
are at least $b_-+1>\eta$.) Since $b/a<1$, all moment bounds above are
bounded by $B^+,E^+$, while
\[
 |h_s^{(0)}|\ge\delta^{s(s-1)},\qquad
 |h_s^{(1)}|\ge\delta^{s^2}.
\]
Therefore a single valid threshold throughout the box is
\[
 C_{\rm box}=\max\left\{1,
 2(k+1)!(k+1)E^+(B^+)^k\delta^{-(k+1)^2}\right\}.
\]
This bound is deliberately conservative, but effective in every original
parameter range and independent of any finite numerical search.

\paragraph{One fully evaluated certificate.}
For $b=1$, $\gamma=1/2$, and $k=1$, the finite recursions give
\[
 C_*(1,1/2,1)=\frac{2\,454\,627}{2}.
\]
Consequently every real $c\ge2\,454\,627/2$ has
$\alpha_1(c),\alpha_2(c)>0$ and $\alpha_3(c)<0$.
This is a rigorous sufficient bound, not a claim of an optimal threshold.
Direct exact arithmetic at $c=1\,227\,314$ independently verifies the
three signs, and the illustrative table in Section~\ref{sec:checks}
shows that successful points can occur much earlier.


\subsection{Existing zero theorems and delayed sign detection}
\label{sec:zeros}
The conclusion here is a specified coefficient sign. Eventual failure
of $\LP$, the subject of Sokal 8(d), already follows from his
Theorem 5 \cite{sokal2022}. A stronger zero-count result is also known.
Cho--Chung--Park \cite{cho2026}, Definition 2.1 and Proposition 5.1(i),
prove that
\[
 \Phi(z)={}_1F_2(a;b,c;-z^2/4)
\]
has infinitely many nonreal zeros if
\[
 b+c>3a+\tfrac12,\qquad
 b-a,c-a\notin\{0,-1,-2,\ldots\}.
\]
For $a=b+\gamma$ with positive nonintegral $\gamma$, these hypotheses
hold whenever $c>2b+3\gamma+\tfrac12$; this also gives $c>a$.
All coefficients of $f_c$ are positive, so $f_c(x)>0$ for $x\ge0$.
Consequently $\Phi$ has no purely imaginary zeros. A nonreal zero of
$\Phi$ is therefore mapped by $z\mapsto-z^2/4$ to a nonreal zero of
$f_c$. This finite-to-one map yields infinitely many nonreal zeros of
$f_c$, and real coefficients arrange them in conjugate pairs.

Combining this established theorem with Theorem~\ref{thm:sokal} gives
the precise comparison
\begin{equation}\label{eq:delayed}
 c>\max\{C_*(b,\gamma,k),\,2b+3\gamma+\tfrac12\}.
\end{equation}
Under \eqref{eq:delayed}, $f_c$ has infinitely many nonreal conjugate
zero pairs, yet its first $2k$ prescribed $S$-fraction coefficients are
positive and its next coefficient is negative. The infinitely-many-zero
assertion is credited to \cite{cho2026}; the finite detection depth is
the assertion proved here. No new zero-count range is claimed.

\section{Prior results and the six research gates}\label{sec:priority}
\subsection{What overlaps with prior mathematics}
Harmonic reflection and low-degree polynomial summation modulo a prime
are established techniques. Meng--Sun \cite{mengsun2015} use related
reflection and harmonic cancellations in a different, single-sum
supercongruence. The direct proof following
Lemma~\ref{lem:antiderivative} makes the family resemblance explicit.
The general derivative statement here is proved in full, but naming
its discrete primitive does not establish a new historical principle.

Z.-H.~Sun's 2022 paper \cite{zhsun2022} gives fourth-power centered
row congruences and a cubic-to-linear moment identity. Those results
substantially overlap with the fourth-power specializations of the
present theorem. In particular, they supply previously available row
estimates for the first two clauses of Sun 36(iii); an elementary
convolution reduction then suffices. These clauses must not be counted
as wholly new vanishing estimates merely because the double-sum
notation is different. The general theorem here allows arbitrary
integer powers $m\ge2$ and odd weights subject to $ms+h<p-1$;
whether this full theorem or its fifth-power applications duplicate a
prior result remains a bibliographic question.

\paragraph{Explicit reduction of the fourth-power clauses.}
For clarity, set $m=4$ in the notation $F_r$ of
Lemma~\ref{lem:row-jet}, and orient $X=-a$, $Y=a-1$ so that their
residues are $s$ and $p-1-s$. In these clauses $s\ge1$ and
$4s+3<p-1$, hence $\langle Y\rangle_p>p/2$; both $X,Y$ are units
and neither has residue $-1$. Theorem 5.1 and equation (5.2) of
\cite{zhsun2022}, printed p.~18, give
$F_1(Y)\equiv0\pmod{p^2}$. Lemma 2.1 of that paper, printed p.~5,
specialized to fourth powers, gives the exact identity
\[
 F_3(z)=-z^2F_1(z)-2z^3\binom{z-1}{p-1}^{4}.
\]
For the parameters in question the final term is zero modulo $p^4$,
so $F_3(Y)\equiv0\pmod{p^2}$ as well. Reflection at the smaller
integer residue gives $F_1(X),F_3(X)\equiv0\pmod p$.
The elementary complete-field sums of degrees $4s$ and $4s+2$ give
$F_0(Y),F_2(Y)\equiv0\pmod p$.
Rectangularization and weight expansion now give
\begin{align*}
 C_{p,4,3}(a)\equiv{}&F_3(X)F_0(Y)+3F_2(X)F_1(Y)\\
 &+3F_1(X)F_2(Y)+F_0(X)F_3(Y)\equiv0\pmod{p^2}.
\end{align*}
Thus the overlap is an explicit deduction, not just a similarity of
titles or subject matter.

The nonalternating fifth-power clause of 36(ii) reduces directly to
a 2008 theorem, as its proof explicitly shows. It fails the requested
test for an independent new result. The present proof of the other
clauses remains valid regardless of that historical overlap.

On the analytic side, the Kummer continued fraction, confluence, and
Hankel extraction formulas are classical \cite{dlmf,sokal2023}.
The useful additional statements here are the application to the
exact subclaim 8(c), the first negative index, and the explicit
parameter thresholds. We do not prove positivity of Sokal's primitive
denominator polynomials for every $c>0$, nor the degree and leading
coefficient formula for his primitive numerator polynomials. In
particular, replacing a reduced numerator by an unreduced Hankel
numerator would not solve 8(b).

\subsection{Audit outcome}
The following table reports the evidence, rather than assigning a
global success label unsupported by that evidence.

\small
\begin{longtable}{@{}>{\raggedright\arraybackslash}p{0.07\linewidth}>{\raggedright\arraybackslash}p{0.22\linewidth}>{\raggedright\arraybackslash}p{0.64\linewidth}@{}}
\toprule
Gate & Requirement & Outcome for this manuscript\\
\midrule
\endhead
G1 & Pre-existing named target & Verified for Sun 36(iii), the literal clauses of 36(ii), and Sokal 8(c), with primary sources and publication dates.\\[4pt]
G2 & Original statement preserved & Exact binomial/Pochhammer identity established before application; original signs, moduli, prime ranges, and Sokal normalization retained. The ambiguous last formula of Sun 36(i) is not counted.\\[4pt]
G3 & Actually open at the time of solution & The targets were published as conjectures in 2019 and 2022. That does not establish their status in November 2026. Some clauses already follow from earlier theorems; current global open status is not certified.\\[4pt]
G4 & Unconditional complete proof & Detailed proofs close all quantifiers of the stated general congruence and Sokal 8(c). Established inputs are named. Finite checks are supplementary; no conjectural corpus result is assumed. External peer review and formal verification have not occurred.\\[4pt]
G5 & Substantive advance, not repackaging & The arbitrary-power degree criterion and finite determinant error bound are substantive statements. Several specializations and all basic tools have antecedents. A historically new methodology is not established, and the known-theorem reductions are explicitly excluded from novelty credit.\\[4pt]
G6 & No prior identical theorem or essential proof & Not certified. Exact-target searches and comparisons with primary papers found important overlaps. Failure to locate an exact duplicate cannot prove priority.\\
\bottomrule
\end{longtable}
\normalsize

Thus this document supplies complete mathematical arguments for the
identified clauses, but does not assert that the user's conjunction
G1--G6 has been satisfied. Nor does it rank these results above every
claim in the supplied corpus. Such a ranking would require validating
those claims and establishing the historical status of the present
ones.

\section{Reproducibility and exact checks}\label{sec:checks}
Every proof above is symbolic and covers its whole stated range.
As an independent transcription check, direct summation of the original
triangle was performed for 39,200 admissible tuples: odd primes
$p\le31$, $2\le m\le8$, $h\in\{1,3,5,7\}$, and every lift
$a\bmod p^2$ satisfying the degree condition. All residues were zero.
For the four literal clauses of 36(iii), every applicable prime
$p\le251$ was tested: the clause counts were $24,25,12,14$.
A separate direct implementation checked 9,234 smaller admissible
tuples and the source clauses through $p=101$, with the same outcome.
These are finite checks of the formulas, not the proof of the theorem.

For the analytic statement, exact rational series inversion, without
Hankel formulas or floating-point sign decisions, gave the following
examples. In this table $b=1$ and $\gamma=k-\tfrac12$.
The listed $c$ values are illustrative successful points, not claimed
to equal the sufficient bound $C_*$.
\begin{center}
\begin{tabular}{@{}rrcc@{}}
\toprule
$k$ & $c$ & Signs of $\alpha_1,\ldots,\alpha_{2k+1}$ & First negative index\\
\midrule
1 & $10^2$ & $++-$ & 3\\
2 & $10^3$ & $++++-$ & 5\\
3 & $10^5$ & $++++++-$ & 7\\
4 & $10^6$ & $++++++++-$ & 9\\
5 & $10^8$ & $++++++++++-$ & 11\\
6 & $10^{10}$ & $++++++++++++-$ & 13\\
\bottomrule
\end{tabular}
\end{center}

Appendix~\ref{app:code} embeds a complete standard-library Python
checker in this single source file. It evaluates the original
triangular sums, extracts continued-fraction coefficients by series
inversion, and computes the explicit threshold. No external data,
bibliography file, image, or separate code file is needed to compile
the manuscript. Compile twice with \texttt{pdflatex}; the first run
creates references and the second resolves them.

\appendix
\section{An exact arithmetic verifier}\label{app:code}
The following Python 3 program uses integer modular arithmetic and
\texttt{fractions.Fraction}. The threshold check at the end uses the
theorem's explicit bound, but verifies the resulting sign by direct
formal inversion. Copy the listing as one program to run it.
\begin{lstlisting}
from fractions import Fraction as F
from math import factorial, prod, ceil


def primes(bound):
    for p in range(2, bound + 1):
        if all(p % d for d in range(2, int(p**0.5) + 1)):
            yield p


def triangle(p, a, m, h):
    mod = p*p
    rows = []
    for z in (a % mod, (1-a) % mod):
        row, u = [1], 1
        for j in range(1, p):
            u = u * (z+j-1) * pow(j, -1, mod) % mod
            row.append(pow(u, m, mod))
        rows.append(row)
    return sum(pow(2*k+1, h, mod) *
               sum(rows[0][j]*rows[1][k-j]
                   for j in range(k+1))
               for k in range(p)) % mod


def exact_s_fraction(b, gamma, c, depth):
    a = b + gamma
    u = [F(1)]
    for j in range(1, depth+2):
        u.append(u[-1]*(a+j-1)*c /
                 ((b+j-1)*(c+j-1)*j))
    q = []
    for j in range(depth+1):
        q.append((j+1)*u[j+1] -
                 sum((u[i]*q[j-i] for i in range(1, j+1)), F(0)))
    g = [b/a * (-1)**j * q[j] for j in range(depth+1)]
    alpha = []
    for unused in range(depth):
        first = g[1]
        assert first != 0
        alpha.append(first/c)
        inverse = [F(1)]
        for j in range(1, len(g)):
            inverse.append(-sum((g[i]*inverse[j-i]
                                  for i in range(1, j+1)), F(0)))
        g = [-inverse[j+1]/first for j in range(len(g)-1)]
    return alpha


def effective_threshold(b, gamma, k):
    a, n = b+gamma, 2*k+1
    u = [F(1)]
    for j in range(1, n+2):
        u.append(u[-1]*(a+j-1)/((b+j-1)*j))
    weight = [F(j*(j-1), 2) for j in range(n+2)]
    q, e = [], []
    for j in range(n+1):
        q.append((j+1)*u[j+1] +
                 sum((u[i]*q[j-i] for i in range(1, j+1)), F(0)))
        e.append((j+1)*u[j+1]*weight[j+1] +
                 sum((u[i]*(e[j-i]+weight[i]*q[j-i])
                      for i in range(1, j+1)), F(0)))
    big_b = max([F(1)] + [b/a*x for x in q])
    big_e = max(b/a*x for x in e)
    beta = [F(0)]  # one-based indices
    for j in range(k+1):
        beta.append((gamma-j)/((b+2*j)*(b+2*j+1)))
        beta.append((a+j+1)/((b+2*j+1)*(b+2*j+2)))
    bound = F(1)
    for s in range(1, k+2):
        h0 = prod((beta[2*j-1]*beta[2*j])**(s-j)
                  for j in range(1, s))
        h1 = h0*prod(beta[2*j-1] for j in range(1, s+1))
        for determinant in (h0, h1):
            assert determinant != 0
            bound = max(bound, 2*factorial(s)*s*big_e *
                        big_b**(s-1)/abs(determinant))
    return bound


def main():
    count = 0
    for p in primes(31):
        if p == 2:
            continue
        for m in range(2, 9):
            for h in (1, 3, 5, 7):
                for a0 in range(p):
                    s = min((-a0) % p, (a0-1) % p)
                    if m*s+h >= p-1:
                        continue
                    for lift in range(p):
                        a = a0 + p*lift
                        assert triangle(p, a, m, h) == 0
                        count += 1
    assert count == 39200
    print('General tuples:', count)
    counts = [0, 0, 0, 0]
    for p in primes(251):
        if p in (2, 5):
            continue
        fifth = pow(5, -1, p*p)
        clauses = [
            (p > 11 and p % 5 in (1, 4), fifth, 4, 3),
            (p > 11 and p % 5 in (2, 3), 2*fifth, 4, 3),
            (p % 5 == 4, fifth, 5, 1),
            (p % 5 == 3, 2*fifth, 5, 1)]
        for i, (active, a, m, h) in enumerate(clauses):
            if active:
                assert triangle(p, a, m, h) == 0
                counts[i] += 1
    assert counts == [24, 25, 12, 14]
    print('Sun 36(iii) clauses:', counts)
    for k, c in enumerate((100, 1000, 10**5, 10**6,
                           10**8, 10**10), start=1):
        b, gamma = F(1), F(2*k-1, 2)
        alpha = exact_s_fraction(b, gamma, F(c), 2*k+1)
        assert all(x > 0 for x in alpha[:-1]) and alpha[-1] < 0
        print('Sokal example:', k, c, '+'*(2*k)+'-')
    # One exact end-to-end check of the proved sufficient bound.
    bound = effective_threshold(F(1), F(1, 2), 1)
    c = F(ceil(bound))
    alpha = exact_s_fraction(F(1), F(1, 2), c, 3)
    assert alpha[0] > 0 and alpha[1] > 0 and alpha[2] < 0
    print('C_*(1, 1/2, 1) =', bound, '; ceil =', c)


if __name__ == '__main__':
    main()
\end{lstlisting}

\section{Source identification and search limitations}
The bibliography distinguishes Zhi-Wei Sun, who posed Conjecture 36,
from Zhi-Hong Sun, who proved the 2022 single-row formulas. These are
different authors. The 2019 journal text, rather than the differently
numbered earlier version of arXiv:0911.5665, supplies the conjecture
number and formulas used in this manuscript.

The literature comparison was conducted through 24 November 2026.
It included the original sources, exact conjecture labels, variations
of the displayed sums, the relevant single-row and double-sum
supercongruence literature, and papers on hypergeometric
Laguerre--P\'olya classification and continued fractions. This was a
targeted research review, not an exhaustive mathematical bibliography.
Search engines may miss papers, alternate notation, unpublished
arguments, or results embedded in more general theorems. Hence no
absence-of-search-result inference is used to certify G3 or G6.

\clearpage
\begin{thebibliography}{10}
\bibitem{sun2019}
Z.-W.~Sun, \emph{Open conjectures on congruences}, Nanjing Univ.
J. Math. Biquarterly \textbf{36} (2019), no.~1, 1--99.
\href{https://doi.org/10.3969/j.issn.0469-5097.2019.01.01}{doi:10.3969/j.issn.0469-5097.2019.01.01}.
Author-uploaded text:
\url{https://www.researchgate.net/publication/336577970_Open_Conjectures_on_Congruences}.
Conjecture 36, printed pp.~31--32.

\bibitem{mengsun2015}
X.-Z.~Meng and Z.-W.~Sun, \emph{Proof of a conjectural supercongruence},
Finite Fields Appl. \textbf{35} (2015), 86--91.
\url{https://arxiv.org/abs/1502.06909}.

\bibitem{zhsun2022}
Z.-H.~Sun, \emph{Supercongruences for sums involving
$\binom ak^m$}, preprint, 2022, arXiv:2204.10132.
\url{https://arxiv.org/abs/2204.10132}.
See Lemma 2.1 and Theorems 2.1 and 5.1.

\bibitem{mccarthyosburn2008}
D.~McCarthy and R.~Osburn, \emph{A $p$-adic analogue of a formula
of Ramanujan}, Arch. Math. \textbf{91} (2008), 492--504.
\href{https://doi.org/10.1007/s00013-008-2828-0}{doi:10.1007/s00013-008-2828-0}.
\url{https://arxiv.org/abs/0708.3307}. Theorem 1.2.

\bibitem{sokal2022}
A.~D.~Sokal, \emph{When does a hypergeometric function ${}_pF_q$
belong to the Laguerre--P\'olya class $LP^+$?}, J. Math. Anal.
Appl. \textbf{515} (2022), 126432.
\href{https://doi.org/10.1016/j.jmaa.2022.126432}{doi:10.1016/j.jmaa.2022.126432}.
\url{https://arxiv.org/abs/2204.01045}. Theorem 5 and Appendix,
Conjecture 8.

\bibitem{sokal2023}
A.~D.~Sokal, \emph{A simple algorithm for expanding a power series
as a continued fraction}, Exposition. Math. \textbf{41} (2023),
245--287.
\href{https://doi.org/10.1016/j.exmath.2022.12.001}{doi:10.1016/j.exmath.2022.12.001}.
\url{https://arxiv.org/abs/2206.15434}. Theorem 9.5(b,c).

\bibitem{dlmf}
NIST Digital Library of Mathematical Functions, \emph{Confluent
hypergeometric functions: continued fractions}, \S13.5,
equations 13.5.1--13.5.2.
\url{https://dlmf.nist.gov/13.5}. Accessed 24 November 2026.

\bibitem{cho2026}
Y.-K.~Cho, S.-Y.~Chung, and Y.~W.~Park,
\emph{A sampling theorem of hypergeometric ${}_1F_2$ functions and
the Laguerre--P\'olya class}, Numer. Algorithms \textbf{102}
(2026), 1713--1741; published online 5 July 2025.
\href{https://doi.org/10.1007/s11075-025-02158-x}{doi:10.1007/s11075-025-02158-x}.
Definition 2.1 and Proposition 5.1(i).

\bibitem{corpus}
QEEDDD, \emph{Verifiable Mathematical Reasoning Corpus}, GitHub
repository, snapshot reviewed 24 November 2026.
\url{https://github.com/QEEDDD/Verifiable-Mathematical-Reasoning-Corpus/tree/52c1e4ae1f82bd3fdc8b809b9a94e2e4c6ef4cb3}.
Relevant inspected entries: [230] \texttt{Catalan\_Hankel\_and\_Sun\_Congruences.tex};
[255] \texttt{Second\_Minima\_and\_Supercongruences.tex};
[166] \texttt{cm\_power\_preservation.tex}.
\end{thebibliography}
\end{document}
